\chapter{2025年北京卷}
\section{单选题}

\begin{problem}
集合 \(M = \{x\mid 2x - 1 > 5\}\)，\(N=\{1,2,3\}\)，则 \(M\cap N=\)（\quad）

\choices
    {\(\{1,2,3\}\)}
    {\(\{2,3\}\)}
    {\(\{3\}\)}
    {\(\emptyset\)}
\end{problem}

\begin{answer}
D
\end{answer}

\begin{solution}
因为 \(M = \{x\mid x > 3\}\)，所以 \(M\cap N=\emptyset\)，故选：D.
\end{solution}

\begin{problem}
已知复数 \(z\) 满足 \(\i\cdot z +2 = 2\i\)，则 \(|z| =\)（\quad）
\choices
    {\(\sqrt{2}\)}
    {\(2\sqrt{2}\)}
    {4}
    {8}
\end{problem}

\begin{answer}
B
\end{answer}

\begin{solution}
由 \(\i\cdot z +2 = 2\i\) 得 \(z=\frac{-2+2\i}{\i}=2+2\i\)，所以 \(|z|=\sqrt{2^2+2^2}=2\sqrt{2}\)，故选：B.
\end{solution}

\begin{problem}
双曲线 \(x^2-4y^2=4\) 的离心率为（\quad）
\choices
    {\(\dfrac{\sqrt{3}}{2}\)}
    {\(\dfrac{\sqrt{5}}{2}\)}
    {5}
    {\(\sqrt{5}\)}
\end{problem}

\begin{answer}
B
\end{answer}

\begin{solution}
设标准式为 \(\dfrac{x^2}{a^2}-\dfrac{y^2}{b^2}=1\)，则 \(a=2\)，\(b=1\)，故 \(c=\sqrt{a^2+b^2}=\sqrt{5}\)，\(e=c/a=\dfrac{\sqrt{5}}{2}\).
\end{solution}

\begin{problem}
为得到函数 \(y=9^x\) 的图象，只需把函数 \(y=3^x\) 的图象上的所有点

\choices
    {\(x\) 变成原来的 \(\dfrac12\) 倍，\(y\) 不变}
    {\(x\) 变成原来的 \(2\) 倍，\(y\) 不变}
    {\(y\) 变成原来的 \(\dfrac13\) 倍，\(x\) 不变}
    {\(y\) 变成原来的 \(3\) 倍，\(x\) 不变}
\end{problem}

\begin{answer}
A
\end{answer}

\begin{solution}
因为 \(y=9^x=3^{2x}\)，所以把函数 \(y=3^x\) 图象上点的横坐标变成原来的 \(\dfrac12\) 倍，纵坐标保持不变后，所得点满足 \(y=9^x\). 故选：A.
\end{solution}

\begin{problem}
已知 \(\{a_n\}\) 是公差不为 0 的等差数列，\(a_1=-2\)，若 \(a_3\)，\(a_4\)，\(a_6\) 成等比数列，则 \(a_{10}=\)（\quad）
\choices
    {\(-20\)}
    {\(-18\)}
    {16}
    {18}
\end{problem}

\begin{answer}
C
\end{answer}

\begin{solution}
设等差数列 \(\{a_n\}\) 的公差为 \(d\)（\(d\neq 0\)）. 因为 \(a_3\)，\(a_4\)，\(a_6\) 成等比数列，所以 \(a_4^2=a_3a_6\)，即 \((-2+3d)^2=(-2+2d)(-2+5d)\). 解得 \(d=2\) 或 \(d=0\). 由 \(d\neq 0\)，得 \(d=2\). 所以 \(a_{10}=a_1+9d=-2+18=16\). 故选：C.
\end{solution}

\begin{problem}
已知 \(a>0\)，\(b>0\)，则（\quad）
\choices
    {\(a^2+b^2>2ab\)}
    {\(\dfrac{1}{a}+\dfrac{1}{b}\ge \dfrac{2}{\sqrt{ab}}\)}
    {\(a+b>\sqrt{ab}\)}
    {\(\dfrac{1}{a}+\dfrac{1}{b}\le \dfrac{2}{\sqrt{ab}}\)}
\end{problem}

\begin{answer}
C
\end{answer}

\begin{solution}
对于 A，当 \(a=b\) 时，\(a^2+b^2=2ab\)，所以 A 错误.

对于 B 与 D，取 \(a=\frac12\)，\(b=\frac14\)，则 \(\frac{1}{a}+\frac{1}{b}=2+4=6\)，\(\frac{2}{\sqrt{ab}}=\frac{2}{\sqrt{\frac12\cdot\frac14}}=4\sqrt2\). 因为 \(6<4\sqrt2\)，所以 B 错误，D 也错误.

对于 C，由基本不等式，\(a+b\ge 2\sqrt{ab}\). 又 \(a\)，\(b>0\)，所以 \(\sqrt{ab}>0\)，从而 \(2\sqrt{ab}>\sqrt{ab}\)，故 \(a+b>\sqrt{ab}\).
所以 C 正确. 故选：C.
\end{solution}

\begin{problem}
已知函数 \(f(x)\) 的定义域为 \(D\)，则“函数 \(f(x)\) 的值域为 \(\R\)”是“对任意 \(M\in\R\)，存在 \(x_0\in D\) 使得 \(|f(x_0)|>M\)”的（\quad）
\choices
    {充分不必要条件}
    {必要不充分条件}
    {充分必要条件}
    {既不充分也不必要条件}
\end{problem}

\begin{answer}
A
\end{answer}

\begin{solution}
若 \(\mathrm{Range}(f)=\R\)，则对任意 \(M\)，取 \(x_1\in D\) 使 \(f(x_1)=|M|+1\)，取 \(x_0=x_1\)，故 \(|f(x_0)|=|M|+1>M\)，所以充分. 反例：\(f(x)=2^x\)，\(D=\R\) 满足对任意 \(M\) 存在 \(x_0\) 使 \(|f(x_0)|>M\)，但值域为 \((0,+\infty)\)，并非 \(\R\). 故选：A.
\end{solution}

\begin{problem}
设函数 \(f(x)=\sin(\omega x)+\cos(\omega x)\)（\(\omega>0\)），若 \(f(x+\pi)=f(x)\) 恒成立，且 \(f(x)\) 在 \(\left[0,\frac{\pi}{4}\right]\) 上有零点，则 \(\omega\) 的最小值为（\quad）
\choices
    {8}
    {6}
    {4}
    {3}
\end{problem}

\begin{answer}
C
\end{answer}

\begin{solution}
因为 \(f(x)=\sqrt{2}\sin\left(\omega x+\frac{\pi}{4}\right)\)，其最小正周期是 \(T=\frac{2\pi}{\omega}\).
由题设 \(f(x+\pi)=f(x)\) 恒成立，可得 \(kT=\pi\)（\(k\in\N^+\)），所以 \(\frac{2\pi}{\omega}=\frac{\pi}{k}\)，即 \(\omega=2k\)（\(k\in\N^+\)）.
又 \(f(x)\) 在 \(\left[0,\frac{\pi}{4}\right]\) 上存在零点. 当 \(x\in\left[0,\frac{\pi}{4}\right]\) 时，\(\omega x+\frac{\pi}{4}\in\left[\frac{\pi}{4},\frac{\pi\omega}{4}+\frac{\pi}{4}\right]\). 要使该区间内能取到 \(\pi\)，必须有 \(\frac{\pi\omega}{4}+\frac{\pi}{4}\ge\pi\)，即 \(\omega\ge3\).
结合 \(\omega=2k\)（\(k\in\N^+\)），可得 \(\omega\) 的最小值为 \(4\). 故选：C.
\end{solution}

\begin{problem}
在一定条件下，某大模型训练 \(N\) 个单位的数据量所需时间 \(T=k\log_2 N\)（单位：小时）. 已知数据量从 \(10^6\) 增加到 \(1.024\times10^9\) 时增加 \(20\) 小时；当数据量从 \(1.024\times10^9\) 增加到 \(4.096\times10^9\) 时，训练时间增加（\quad）
\choices
    {2}
    {4}
    {20}
    {40}
\end{problem}

\begin{answer}
B
\end{answer}

\begin{solution}
设 \(T_1=T(10^6)\)，\(T_2=T(1.024\times10^9)\)，\(T_3=T(4.096\times10^9)\). 因为 \(1.024\times10^9=1024\times10^6=2^{10}\times10^6\)，所以
\[
T_2-T_1=k\log_2 \frac{1.024\times10^9}{10^6}=k\log_2 2^{10}=10k=20
\]
故 \(k=2\). 又因为 \(4.096\times10^9=4\times1.024\times10^9=2^2\times1.024\times10^9\)，所以
\[
T_3-T_2=k\log_2 \frac{4.096\times10^9}{1.024\times10^9}=k\log_2 4=2k=4
\]
故选：B.
\end{solution}

\begin{problem}
已知平面直角坐标系 \(xOy\) 中，\(|\overrightarrow{OA}|=|\overrightarrow{OB}|=\sqrt2\)，\(|\overrightarrow{AB}|=2\)，设 \(C(3,4)\)，则 \(|2\overrightarrow{CA}+\overrightarrow{AB}|\) 的取值范围是（\quad）
\choices
    {\([6,14]\)}
    {\([6,12]\)}
    {\([8,14]\)}
    {\([8,12]\)}
\end{problem}

\begin{answer}
D
\end{answer}

\begin{solution}
由 \(|\overrightarrow{OA}|=|\overrightarrow{OB}|=\sqrt2\)，\(\overrightarrow{OA}\cdot\overrightarrow{OB}=0\). 设 \(\overrightarrow{OA}+\overrightarrow{AB}=\overrightarrow{OB}\). 展开 \(|2\overrightarrow{CA}+\overrightarrow{AB}|^2 =|2(\overrightarrow{OA}-\overrightarrow{OC})+\overrightarrow{AB}|^2\)，并代入 \(|\overrightarrow{OC}|=\sqrt{25}\)，\(\overrightarrow{OC}\) 与 \(\overrightarrow{OA}+\overrightarrow{OB}\) 的内积介于 \(\pm10\) 之间，可得 \(|2\overrightarrow{CA}+\overrightarrow{AB}|^2\in[64,144]\)，故区间为 \([8,12]\).
\end{solution}
\section{填空题}

\begin{problem}
抛物线 \(y^2=2px\)（\(p>0\)）的顶点到焦点的距离为 3，则 \(p=\)\fillinblank{}.
\end{problem}

\begin{answer}
6
\end{answer}

\begin{solution}
抛物线 \(y^2=2px\) 中，顶点到焦点距离为 \(\dfrac p2\)，故 \(\dfrac p2=3\)，则 \(p=6\).
\end{solution}

\begin{problem}
已知 \((1-2x)^4=a_0-2a_1x+4a_2x^2-8a_3x^3+16a_4x^4\)，则 \(a_0=\)；\(a_1+a_2+a_3+a_4=\)\fillinblank{}.
\end{problem}

\begin{answer}
\(1\)；\(15\)
\end{answer}

\begin{solution}
令 \(x=0\) 得 \(a_0=1\). 置 \(t=-2x\) 得 \((1+t)^4=a_0+a_1t+a_2t^2+a_3t^3+a_4t^4\). 令 \(t=1\) 得 \(a_0+a_1+a_2+a_3+a_4=16\)，故 \(a_1+a_2+a_3+a_4=15\).
\end{solution}

\begin{problem}
已知 \(\alpha\)，\(\beta\in[0,2\pi]\)，且 \(\sin(\alpha+\beta)=\sin(\alpha-\beta)\)，\(\cos(\alpha+\beta)\ne\cos(\alpha-\beta)\)，写出满足条件的一组 \(\alpha=\fillinblank{}\)，\(\beta=\fillinblank{}\).
\end{problem}

\begin{answer}
\(\alpha=\frac{\pi}{2}\)，\(\beta=\frac{\pi}{6}\)（答案不唯一）
\end{answer}

\begin{solution}
由 \(\sin(\alpha+\beta)=\sin(\alpha-\beta)\) 有两种可能，其中应排除 \(\cos(\alpha+\beta)=\cos(\alpha-\beta)\) 的情形. 取 \((\alpha,\beta)=\left(\frac{\pi}{2},\frac{\pi}{6}\right)\)，则
\(\sin(\alpha+\beta)=\sin\frac{2\pi}{3}=\frac{\sqrt3}{2}\)，\(\sin(\alpha-\beta)=\sin\frac{\pi}{3}=\frac{\sqrt3}{2}\)，且 \(\cos(\alpha+\beta)=\cos\frac{2\pi}{3}=-\frac12\ne\frac12=\cos\frac{\pi}{3}=\cos(\alpha-\beta)\).
故这组 \(\alpha\)，\(\beta\) 满足题设.
\end{solution}

\begin{problem}
某科技兴趣小组通过 3D 打印机打印的一个零件可以抽象为如图所示的多面体，其中 \(ABCDEF\) 是一个平行六边形，平面 \(ARF\perp\) 平面 \(ABC\)，平面 \(TCD\perp\) 平面 \(ABC\)，\(AB\perp BC\)，\(AB\parallel RS\parallel EF\parallel CD\)，\(AF\parallel ST\parallel BC\parallel ED\). 若 \(AB=BC=8\)，\(AF=CD=4\)，\(AR=RF=TC=TD=\dfrac52\)，则该多面体的体积为（\quad）

\texfigure[max width=0.48\linewidth]{img_repaint/2025/beijing/q14_fig1.tex}

\choices
    {60}
    {90}
    {120}
    {150}
\end{problem}

\begin{answer}
A
\end{answer}

\begin{solution}
先证明一个结论：若平面 \(\alpha\perp\gamma\)，平面 \(\beta\perp\gamma\)，且 \(\alpha\cap\beta=l\)，则 \(l\perp\gamma\). 证明如下：设 \(\alpha\cap\gamma=a\)，\(\beta\cap\gamma=b\). 在平面 \(\gamma\) 内取一点 \(O\)，且 \(O\notin a\)，\(O\notin b\). 在平面 \(\gamma\) 内过 \(O\) 作直线 \(m\perp a\)，再作直线 \(n\perp b\). 因为 \(\alpha\perp\gamma\) 且 \(m\subset\gamma\)，所以 \(m\perp\alpha\)，进而 \(m\perp l\). 又因为 \(\beta\perp\gamma\) 且 \(n\subset\gamma\)，所以 \(n\perp\beta\)，进而 \(n\perp l\). 又 \(m\cap n=O\)，且 \(m\)，\(n\subset\gamma\)，所以 \(l\perp\gamma\). 下面回到原题. 连接 \(BE\). 因为 \(AB\perp BC\) 且 \(AF\parallel BC\)，所以 \(AF\perp AB\). 又 \(BC\perp AB\) 且 \(CD\parallel AB\)，所以 \(BC\perp CD\). 再由 \(EF\parallel CD\)，\(ED\parallel BC\)，得 \(EF\perp ED\). 因为 \(AB\perp AF\)，\(BC\perp CD\)，且 \(AB=BC=8\)，\(AF=CD=4\)，所以直角梯形 \(ABEF\) 与直角梯形 \(CBDE\) 的两条直角边分别对应相等，从而这两个直角梯形全等. 因此 \(\angle BEF=\angle BED=45^\circ\). 在直角梯形 \(ABEF\) 中，过 \(B\) 作 \(BT\perp EF\) 于 \(T\)，则四边形 \(ABTF\) 为矩形，且 \(\triangle BTE\) 为等腰直角三角形，因此
\[EF=FT+TE=AB+BT=AB+AF=12\]
由平面 \(RAF\perp\) 平面 \(ABEF\)，且两平面交于 \(AF\)，又 \(AB\perp AF\) 且 \(AB\subset\) 平面 \(ABEF\)，所以 \(AB\perp\) 平面 \(RAF\). 取 \(AF\) 的中点为 \(M\)，\(BE\) 的中点为 \(U\)，\(CD\) 的中点为 \(V\)，连接 \(RM\)，\(MU\)，\(SU\)，\(UV\). 由中点性质得 \(MU\parallel RS\). 又 \(RM\subset\) 平面 \(RAF\)，且 \(AB\perp\) 平面 \(RAF\)，故 \(RM\perp AB\). 由于 \(AB\subset\) 平面 \(ABEF\)，再由 \(RM\parallel SU\) 可得 \(SU\perp\) 平面 \(ABEF\). 于是四边形 \(RMUS\) 为平行四边形，从而
\[
RS=MU=\frac12\left( BE+AF \right)=\frac12\left( 12+8 \right)=10
\]
在平面 \(ABHR\) 中过 \(B\) 作 \(BH\parallel AR\)，交 \(RH\) 于 \(H\)，连接 \(HT\). 则四边形 \(ABHR\) 为平行四边形，故 \(RH\parallel AB\) 且 \(RH=AB\). 又 \(FT=AB\)，所以 \(RH\parallel FT\) 且 \(RH=FT\)，从而四边形 \(RFTH\) 为平行四边形. 因为 \(BH\perp AB\)，\(BT\perp AB\)，且 \(BH\)，\(BT\subset\) 平面 \(BHT\)，所以 \(AB\perp\) 平面 \(BHT\). 于是平面 \(ARF\parallel\) 平面 \(BHT\). 再由 \(AR=BH\)，\(RF=HT\)，\(AF=BT\)，得 \(\triangle ARF\cong\triangle BHT\)，故几何体 \(ARF-BHT\) 为直棱柱. 在 \(\triangle ARF\) 中，\(S_{\triangle ARF}=\frac12\cdot 4\cdot \sqrt{\left(\frac52\right)^2-2^2}=3\)，所以 \(V_{ARF-BHT}=8\times 3=24\). 又因为 \(AB\parallel EF\)，所以 \(EF\perp\) 平面 \(ARF\). 而 \(EF\subset\) 平面 \(RSEF\)，故平面 \(ARF\perp\) 平面 \(RSEF\). 在平面 \(ARF\) 内过 \(A\) 作 \(AG\perp RF\)，垂足为 \(G\)，则 \(AG\perp\) 平面 \(RSEF\). 由 \(\frac12\cdot AG\cdot RF=3\) 得 \(AG=\frac{12}{5}\). 因此 \(V_{B-HTES}=\frac13\cdot \frac{12}{5}\cdot \frac12(2+4)\cdot \frac52=6\). 由图形的对称性，整个几何体的体积为
\(
V=2(24+6)=60
\)
故选：A.
\end{solution}
\section{解答题}

\begin{problem}
关于定义域为 \(\R\) 的函数 \(f(x)\)，以下说法正确的有（\quad）

\choices
    {存在在 \(\R\) 上单调递增的函数 \(f(x)\)，使 \(f(x)+f(2x)=-x\) 恒成立}
    {存在在 \(\R\) 上单调递减的函数 \(f(x)\)，使 \(f(x)+f(2x)=-x\) 恒成立}
    {存在且有无穷多个函数 \(f(x)\)，满足 \(f(x)+f(-x)=\cos x\)}
    {存在且有无穷多个函数 \(f(x)\)，满足 \(f(x)-f(-x)=\cos x\)}
\end{problem}

\begin{solution}
若存在在 \(\R\) 上单调递增的函数 \(f(x)\)，使 \(f(x)+f(2x)=-x\) 恒成立，则令 \(x=0\) 得 \(f(0)+f(0)=0\)，所以 \(f(0)=0\).
当 \(x>0\) 时，由单调递增性得
\(
f(4x)>f(2x)>f(x)>f(0)=0
\)
故 \(f(4x)+f(2x)>f(2x)+f(x)\)，即 \(-2x>-x\)，
矛盾，故错误.

取 \(f(x)=-\dfrac{x}{3}\)，则成立，正确.

取 \(f(x)=\dfrac{1}{2}\cos x+mx\)，\(m\in\R\)，则成立且有无穷多个，正确.

若 \(f(x)-f(-x)=\cos x\)，代 \(x=0\) 得 \(0=\cos 0=1\)，矛盾. 故答案为：BC.
\end{solution}

\begin{problem}
在 \(\triangle ABC\) 中，\(\cos A=-\dfrac13\)，\(a\sin C=4\sqrt2\).

\begin{enumerate}
\item 求 \(c\)；

\item 在以下三个条件中选择一个作为已知，使得 \(\triangle ABC\) 存在，求 \(BC\) 的高.
\end{enumerate}

\(\text{\circled{1}} \, a=6\)，\(\text{\circled{2}} \, b\sin C= \frac{10\sqrt2}{3}\)，\(\text{\circled{3}} \, S_{\triangle ABC}=10\sqrt2\)
\end{problem}

\begin{solution}
\begin{enumerate}
\item \(\sin A=\sqrt{1-\cos^2A}=\dfrac{2\sqrt2}{3}\). 由正弦定理得 \(a\sin C=c\sin A=4\sqrt2\)，从而 \(c=6\).

\item 如图，若 \(\triangle ABC\) 存在，设其 \(BC\) 边上的高为 \(AD\).
\solutionfigure{\texfigure[max width=0.54\linewidth]{img_repaint/2025/beijing/q16_solution_fig1.tex}}
若选 \circled{1}，则 \(a=6=c\)，从而 \(A=C\). 又 \(\cos A=-\dfrac13<0\)，故 \(A\) 为钝角，于是 \(C\) 也为钝角，矛盾，所以 \(\triangle ABC\) 不存在.

取 \circled{2}. 由正弦定理，\(b\sin C=c\sin B\). 由 \(b\sin C=\dfrac{10\sqrt2}{3}\) 及 \(c=6\) 得 \(6\sin B=\frac{10\sqrt2}{3}\)，从而 \(\sin B=\frac{5\sqrt2}{9}\). 因为 \(A\) 已为钝角，所以 \(B\) 只能为锐角，故 \(\cos B=\sqrt{1-\sin^2 B}=\frac{\sqrt{31}}{9}\).
设 \(AD\) 为 \(BC\) 边上的高，则
\[
AD=c\sin B=6\times\frac{5\sqrt2}{9}=\frac{10\sqrt2}{3}
\]
又因为 \(\cos\angle DAB=\sin B\)，\(\sin\angle DAB=\cos B\)，且 \(\cos A=-\frac13\)，\(\sin A=\frac{2\sqrt2}{3}\)，
所以 \(\cos\angle CAD=\cos\left( A-\angle DAB \right)\)，\(\sin\angle CAD=\sqrt{1-\cos^2\angle CAD}\) 都可以唯一确定，从而 \(CA\)，\(CD\) 也可以唯一确定. 这表明此时三角形 \(ABC\) 存在，且
\(AD=\frac{10\sqrt2}{3}\).

若选 \(\circled{3}\)，由面积公式 \(S_{\triangle ABC}=\frac12 bc\sin A\)，结合 \(S_{\triangle ABC}=10\sqrt2\)，\(c=6\)，\(\sin A=\frac{2\sqrt2}{3}\) 得 \(b=5\). 再由余弦定理
\[
a^2=b^2+c^2-2bc\cos A=25+36+20=81
\]
故 \(a=9\). 此时 \(A\)，\(B\) 与 \(C\) 可唯一确定，三角形存在. 高为
\(AD=\frac{2S_{\triangle ABC}}a=\frac{20\sqrt2}{9}\).
因此可选的条件为 \(\circled{2}\)，\(\circled{3}\)，对应的 \(BC\) 高分别为 \(\frac{10\sqrt2}{3}\)，\(\frac{20\sqrt2}{9}\).
\end{enumerate}
\end{solution}

\begin{problem}
四棱锥 \(P-ABCD\)，\(\triangle ACD\) 与 \(\triangle ABC\) 为等腰直角三角形，\(\angle ADC=90^\circ\)，\(\angle BAC=90^\circ\)，\(E\) 为 \(BC\) 中点.

\begin{enumerate}
\item \(F\) 为 \(PD\) 中点，\(G\) 为 \(PE\) 中点，证明 \(FG\parallel\) 平面 \(PAB\)；

\item 若 \(PA\perp\) 平面 \(ABCD\)，\(\ PA=AC\)，求 \(AB\) 与平面 \(PCD\) 所成角的正弦值.
\end{enumerate}
\end{problem}

\begin{solution}
\begin{enumerate}
\item 取 \(PA\) 的中点 \(N\)，\(PB\) 的中点 \(M\)，连接 \(FN\)，\(MN\). 不妨设 \(AD=CD=2\)，则由 \(\triangle ACD\) 与 \(\triangle ABC\) 为等腰直角三角形可得 \(AC=AB=2\sqrt2\)，\(BC=4\). 又 \(F\)，\(E\) 分别为 \(PD\)，\(BC\) 的中点，所以 \(FN=\frac12 AD=1\)，\(GM=\frac12 BE=1\)，
且 \(FN\parallel AD\)，\(GM\parallel BC\). 又因为 \(\angle DAC=\angle ACB=45^\circ\)，且 \(AC\) 为截线，所以 \(AD\parallel BC\)，从而 \(FN\parallel GM\). 故四边形 \(FGMN\) 为平行四边形，于是
\(FG\parallel MN\). 因为 \(MN\subset\) 平面 \(PAB\)，且 \(FG\not\subset\) 平面 \(PAB\)，所以 \(FG\parallel \text{平面 }PAB\).

\item 取\(A(0,0,0)\)，\(B(0,2\sqrt2,0)\)，\(C(2\sqrt2,0,0)\)，\(D(\sqrt2,-\sqrt2,0)\)，\(P(0,0,2\sqrt2)\).
\solutionfigure{\texfigure[max width=0.46\linewidth]{img_repaint/2025/beijing/q17_solution_fig1.tex}}
则 \(\overrightarrow{AB}=(0,2\sqrt2,0)\)，\(\overrightarrow{PC}=(-2\sqrt2,0,2\sqrt2)\). 设平面 \(PCD\) 法向量为 \(\bs{n}=(1,-1,1)\)，可由 \(\overrightarrow{CD}\cdot\bs{n}=0\)，\(\overrightarrow{CP}\cdot\bs{n}=0\) 得到. 设 \(\theta\) 为 \(AB\) 与平面 \(PCD\) 的夹角，则
\(\sin\theta=|\cos\angle(\overrightarrow{AB},\bs{n})|=\frac{|(0,2\sqrt2,0)\cdot\bs{n}|}{|\overrightarrow{AB}|\cdot|\bs{n}|}=\frac{\sqrt3}{3}\).
故答案 \(\dfrac{\sqrt3}{3}\).
\end{enumerate}
\end{solution}

\begin{problem}
有一道选择题考查某知识点. 甲，乙两校各随机抽取 \(100\) 人，甲校有 \(80\) 人答对，乙校有 \(75\) 人答对，用频率估计概率.

\begin{enumerate}
    \item 从甲校随机抽取 1 人，求这个人做对该题目的概率；
    \item 从甲乙两校各随机抽取 1 人，设 \(X\) 为做对人数，求 \(P(X=1)\) 与 \(E(X)\)；
    \item 若甲校同学掌握该知识点，则有 \(100\%\) 的概率做对该题目；乙校同学掌握该知识点，则有 \(85\%\) 的概率做对该题目；未掌握该知识点的同学都是从四个选项里面随机选择一个. 设甲校学生掌握该知识点的概率为 \(p_1\)，乙校学生掌握该知识点的概率为 \(p_2\)，比较 \(p_1\) 与 \(p_2\) 的大小.
\end{enumerate}
\end{problem}

\begin{solution}
\begin{enumerate}
\item \(P(\text{甲校随机一人答对})=80/100=0.8\).

\item 设甲校一人答对事件为 \(A\)，乙校一人答对事件为 \(B\)，则
\(P(X=1)=P(A\cap\bar B)+P(\bar A\cap B)=0.8\times0.25+0.2\times0.75=0.35\).
且 \(P(X=1)=0.35\)，\(P(X=2)=0.6\)，\(P(X=0)=0.05\)，故 \(E(X)=0.35+2\times0.6=1.55\).

\item 由 \(0.8=p_1+\frac14(1-p_1)\) 得 \(p_1=\frac{11}{15}\)，由 \(0.75=0.85p_2+\frac14(1-p_2)\) 得 \(p_2=\frac56\)，
故 \(p_1<p_2\).
\end{enumerate}
\end{solution}

\begin{problem}
已知椭圆 \(\dfrac{x^2}{a^2}+\dfrac{y^2}{b^2}=1\)（\(a>b>0\)），其离心率为 \(\dfrac{\sqrt{2}}{2}\)，椭圆上一点到两焦点距离之和为 \(4\).

\begin{enumerate}
\item 求椭圆方程；
\item 设 \(O=(0,0)\)，\(M(x_0,y_0)\)（\(x_0\ne0\)）为椭圆上一点，直线 \(x_0x+2y_0y-4=0\) 与直线 \(y=2\)，\(y=-2\) 交于 \(A\)，\(B\). \(S_1\)，\(S_2\) 分别为 \(\triangle OAM\) 与 \(\triangle OBM\) 的面积，比较 \(\dfrac{S_1}{S_2}\) 与 \(\dfrac{|OA|}{|OB|}\) 的大小.
\end{enumerate}
\end{problem}

\begin{solution}
\begin{enumerate}
\item 椭圆上任意一点到两焦点距离之和恒为 \(2a\). 题设给出这个和等于 \(4\)，所以 \(2a=4\)，从而 \(a=2\). 又 \(e=\frac{c}{a}=\frac{\sqrt{2}}{2}\)，所以 \(c=\sqrt{2}\)，从而 \(b^2=a^2-c^2=2\). 故椭圆方程为
\(
\frac{x^2}{4}+\frac{y^2}{2}=1
\)

\item 联立 \(\begin{cases}
x_0x+2y_0y-4=0,\\[0.5ex]
\dfrac{x^2}{4}+\dfrac{y^2}{2}=1
\end{cases}\) 可得
\(\left(\frac{4-2y_0y}{x_0}\right)^2+2y^2=4\).
结合 \( \frac{x_0^2}{4}+\frac{y_0^2}{2}=1 \) 可化为 \( (y-y_0)^2=0 \)，所以直线 \(x_0x+2y_0y-4=0\) 与椭圆相切，\(M\) 为切点. 设 \(A(x_1,y_1)\)，\(B(x_2,y_2)\)，且 \(x_2<x_0<x_1\). 由 \(y=2\) 与 \(y=-2\) 分别代入直线方程，得 \(x_1=\frac{4-4y_0}{x_0}\)，\(y_1=2\)；\(x_2=\frac{4+4y_0}{x_0}\)，\(y_2=-2\). 于是
\[
\frac{S_1}{S_2}=\frac{|AM|}{|BM|}=\frac{|x_1-x_0|}{|x_2-x_0|}=\frac{2-y_0}{2+y_0}
\]
又
\(
\frac{|OA|}{|OB|}=\frac{\sqrt{\left(\frac{4-4y_0}{x_0}\right)^2+4}}{\sqrt{\left(\frac{4+4y_0}{x_0}\right)^2+4}}=\frac{2-y_0}{2+y_0}
\)
故 \(\frac{S_1}{S_2}=\frac{|OA|}{|OB|}\).

法二：不妨设 \(A(x_1,y_1)\)，\(B(x_2,y_2)\). 当 \(x_1=x_2\) 时，由对称性可知，\(\frac{S_1}{S_2}=\frac{|OA|}{|OB|}\). 故设 \(x_2<x_0<x_1\). 由 \(\begin{cases}
x_0x+2y_0y-4=0,\\
y=2
\end{cases}\) 解得
\(x_1=\frac{4-4y_0}{x_0}\)，\(y_1=2\).
由 \(\begin{cases}
x_0x+2y_0y-4=0,\\
y=-2
\end{cases}\) 解得
\(x_2=\frac{4+4y_0}{x_0}\)，\(y_2=-2\). 于是 \(k_{OA}=\frac{y_1}{x_1}=\frac{2x_0}{4-4y_0}=\frac{x_0}{2-2y_0}\)，\(k_{OB}=\frac{y_2}{x_2}=\frac{-2x_0}{4+4y_0}=-\frac{x_0}{2+2y_0}\)，\(k_{OM}=\frac{y_0}{x_0}\).
又因为
\(\frac{x_0^2}{4}+\frac{y_0^2}{2}=1\)，所以 \(x_0^2+2y_0^2=4\). 从而
\begin{align*}
\tan\angle AOM&=\frac{k_{OA}-k_{OM}}{1+k_{OA}k_{OM}}
=\frac{\dfrac{x_0}{2-2y_0}-\dfrac{y_0}{x_0}}{1+\dfrac{x_0}{2-2y_0}\cdot\dfrac{y_0}{x_0}} \\
&=-\frac{x_0^2+2y_0^2-2y_0}{x_0(y_0-2)}=-\frac{4-2y_0}{x_0(y_0-2)}=\frac{2}{x_0}, \\
\tan\angle BOM&=\frac{k_{OM}-k_{OB}}{1+k_{OM}k_{OB}}
=\frac{\dfrac{y_0}{x_0}+\dfrac{x_0}{2+2y_0}}{1-\dfrac{y_0}{x_0}\cdot\dfrac{x_0}{2+2y_0}} \\
&=\frac{x_0^2+2y_0^2+2y_0}{x_0(y_0+2)}=\frac{4+2y_0}{x_0(y_0+2)}=\frac{2}{x_0}
\end{align*}
于是 \(\tan\angle AOM=\tan\angle BOM\)，即 \(\angle AOM=\angle BOM\). 所以
\(
\frac{S_1}{S_2}=\frac{|OA|\cdot |OM|\sin\angle AOM}{|OB|\cdot |OM|\sin\angle BOM}=\frac{|OA|}{|OB|}
\)
\end{enumerate}
\end{solution}

\begin{problem}
函数 \(f:(-1,+\infty)\to\R\)，\(f(0)=0\)，\(f'(x)=\dfrac{\ln(1+x)}{1+x}\)，\(l_1\) 为点 \(A(a,f(a))\)（\(a\ne0\)）处切线.

\begin{enumerate}
    \item \(f'(x)\) 的最大值；
    \item 证明：当 \(-1<a<0\) 时，除点 \(A\) 外曲线 \(y=f(x)\) 在 \(l_1\) 上方；
    \item 若 \(a>0\) 时，过 \(A\) 且与 \(l_1\) 垂直的 \(l_2\) 与 \(l_1\) 分别交 \(x\) 轴于 \(x_1\)，\(x_2\)，求 \(\dfrac{2a-x_1-x_2}{x_2-x_1}\) 的范围.
\end{enumerate}
\end{problem}

\begin{solution}
\begin{enumerate}
\item 设 \(g(x)=f'(x)=\dfrac{\ln(1+x)}{1+x}\)，则
\(g'(x)=\frac{1-\ln(1+x)}{(1+x)^2}\).
由 \(g'(x)=0\) 得 \(x=\e-1\). 当 \(x\in(-1,\e-1)\) 时，\(g'(x)>0\)；当 \(x\in(\e-1,+\infty)\) 时，\(g'(x)<0\). 所以
\(f'(x)_{\max}=f'(\e-1)=\frac{1}{\e}\).

\item 因为
\(f'(a)=\frac{\ln(1+a)}{1+a}\)，
所以切线 \(l_1\) 的方程为
\(y-f(a)=\frac{\ln(1+a)}{1+a}(x-a)\)，
即
\(y=\frac{\ln(1+a)}{1+a}(x-a)+f(a)\).
设
\(h(x)=f(x)-\left[\frac{\ln(1+a)}{1+a}(x-a)+f(a)\right]\)，
则
\(h(a)=0\)，\(h'(x)=f'(x)-f'(a)=\frac{\ln(1+x)}{1+x}-\frac{\ln(1+a)}{1+a}\).
由（1）知 \(f'(x)\) 在 \((-1,\e-1)\) 上单调递增，而 \(-1<a<0<\e-1\). 因此：当 \(-1<x<a\) 时，\(h'(x)<0\)，所以 \(h(x)\) 在 \((-1,a)\) 上单调递减；

当 \(a<x<0\) 时，\(h'(x)>0\)，所以 \(h(x)\) 在 \((a,0)\) 上单调递增；

又因为 \(f'(a)<f'(0)=0\)，且当 \(x\ge 0\) 时，
\(f'(x)=\frac{\ln(1+x)}{1+x}\ge 0>f'(a)\)，
所以 \(h'(x)>0\)，即 \(h(x)\) 在 \([0,+\infty)\) 上单调递增.

由于 \(h(x)\) 在 \((-1,a)\) 上递减，在 \((a,0)\) 上递增，在 \([0,+\infty)\) 上也递增，所以 \(x=a\) 是整个定义域内的唯一极小点，且最小值为 \(h(a)=0\).
因此当 \(x\ne a\) 时，\(h(x)>0\). 也就是说，除点 \(A\) 外，曲线 \(y=f(x)\) 在切线 \(l_1\) 上方.

\item 由
\(f'(x)=\frac{\ln(1+x)}{1+x}\)
可设
\(f(x)=\frac{\ln^2(1+x)}{2}+C\).
又 \(f(0)=0\)，所以 \(C=0\)，即
\(f(x)=\frac{\ln^2(1+x)}{2}\).
于是
\(y=\frac{\ln(1+a)}{1+a}(x-a)+\frac{\ln^2(1+a)}{2}\)，
\(y=-\frac{1+a}{\ln(1+a)}(x-a)+\frac{\ln^2(1+a)}{2}\)
分别是 \(l_1\)，\(l_2\) 的方程.

令 \(y=0\)，得
\(x_1=a-\frac{(1+a)\ln(1+a)}{2}\)，\(x_2=a+\frac{\ln^3(1+a)}{2(1+a)}\).
因此
\(
\frac{2a-x_1-x_2}{x_2-x_1}
=\frac{1-\left(\dfrac{\ln(1+a)}{1+a}\right)^2}{1+\left(\dfrac{\ln(1+a)}{1+a}\right)^2}
=\frac{1-g^2(a)}{1+g^2(a)}
\)
由（1）知，当 \(a>0\) 时，
\(g(a)\in\left(0,\frac{1}{\e}\right]\)，
所以
\(g^2(a)\in\left(0,\frac{1}{\e^2}\right]\).
从而 \(\frac{2a-x_1-x_2}{x_2-x_1}\in\left[\frac{\e^2-1}{\e^2+1},1\right)\).
\end{enumerate}
\end{solution}

\begin{problem}
设 \(M=\{(x_i,y_i)\mid x_i,y_i\in A\}\)，\(A=\{1,2,3,4,5,6,7,8\}\). 从 \(M\) 中选 \(n\) 个有序数对构成一列 \((x_1,y_1)\)，\(\dots\)，\((x_n,y_n)\). 若相邻两项满足 \(|x_{i+1}-x_i|=3\text{\ 或\ }4\)，\(|y_{i+1}-y_i|=4\text{\ 或\ }3\)，\(i=1\)，\(\dots\)，\(n-1\)，则该列称为 \(k\) 列.

\begin{enumerate}
    \item 若 \(k\) 列的第一项为 \((3,3)\)，求第二项.
    \item 若 \(\tau\) 为 \(k\) 列，且 \(i\) 为奇数时 \(x_i\in\{1,2,7,8\}\)，\(i\) 为偶数时 \(x_i\in\{3,4,5,6\}\)，判断 \((3,2)\) 与 \((4,4)\) 能否同时在 \(\tau\) 中，并说明.
    \item 证明 \(M\) 中所有元素都不构成 \(k\) 列.
\end{enumerate}
\end{problem}

\begin{solution}
\begin{enumerate}
\item 由题意，相邻两项满足
\(x_{i+1}=x_i\pm 3\)，\(\ y_{i+1}=y_i\pm 4\) 或 \(x_{i+1}=x_i\pm 4\)，\(\ y_{i+1}=y_i\pm 3\).
若第一项为 \((3,3)\)，则当横坐标差为 \(3\) 时，下一项横坐标只能取 \(6\)；当横坐标差为 \(4\) 时，下一项横坐标只能取 \(7\). 对应地，纵坐标分别取 \(7\) 或 \(6\). 所以第二项只能是
\((6,7)\) 或 \((7,6)\).

\item 假设 \((3,2)\) 与 \((4,4)\) 同时在 \(\tau\) 中. 由于 \(k\) 列倒序后仍是 \(k\) 列，不妨设 \((3,2)\) 在 \((4,4)\) 之前. 在 \(k\) 列中，相邻两项的横纵坐标之和的奇偶性总是相反的，所以从 \((3,2)\) 到 \((4,4)\) 需要经过奇数步. 另一方面，题设规定奇数项的横坐标属于 \(\{1,2,7,8\}\)，偶数项的横坐标属于 \(\{3,4,5,6\}\). 而 \((3,2)\) 与 \((4,4)\) 的横坐标都属于 \(\{3,4,5,6\}\)，所以它们在数列中的位置同为偶数位，从前者走到后者必经过偶数步. 奇数步与偶数步矛盾，因此 \((3,2)\) 与 \((4,4)\) 不能同时在 \(\tau\) 中.

\item 法一：若 \(M\) 中所有元素都构成 \(k\) 列，考虑 \(k\) 列中形如 \((x_i,y_i)\) 且 \(x_i\)，\(y_i\in\{1,2,7,8\}\) 的项，这样的项共有 \(16\) 个. 由题意知其下一项 \((x_{i+1},y_{i+1})\) 满足 \(x_{i+1}\)，\(y_{i+1}\in\{3,4,5,6\}\). 但是 \((x_{i+1},y_{i+1})\ne (3,3)\)，\((6,3)\)，\((3,6)\)，\((6,6)\)，因为 \(6\) 只能由 \(2\) 来，\(3\) 只能由 \(7\) 来，横，纵坐标不能同时相差 \(4\). 这样下一项只能有 \(12\) 个点，即对于 \(16\) 个 \((x_i,y_i)\)，只有 \(12\) 个 \((x_{i+1},y_{i+1})\) 与之对应，矛盾.

法二：反设 \(M\) 中所有元素都能构成一个 \(k\) 列，则该列长度 \(n=64\). 设 \(T_1=\{(x,y)\mid x\in\{1,2,7,8\},\ y\in A\}\)，\(T_2=\{(x,y)\mid x\in\{3,4,5,6\},\ y\in A\}\). 则 \(T_1\) 与 \(T_2\) 都有 \(32\) 个元素，且 \(T_1\) 中元素的相邻项必在 \(T_2\) 中. 如果至少存在两对相邻项都属于 \(T_2\)，那么属于 \(T_2\) 的项数会多于属于 \(T_1\) 的项数，这与二者都恰有 \(32\) 个元素矛盾. 因此至多存在一对相邻项同时属于 \(T_2\).

于是必存在 \(m\in\{0,1,\dots,32\}\)，使得
\(
\begin{cases}
(x_{2k-1},y_{2k-1})\in T_1,\ (x_{2k},y_{2k})\in T_2,& 1\le k\le m,\\
(x_{2k-1},y_{2k-1})\in T_2,\ (x_{2k},y_{2k})\in T_1,& m+1\le k\le 32.
\end{cases}
\)

又因为相邻两项的横纵坐标之和奇偶性相反，所以 \(T_1\) 中横纵坐标之和为奇数的点与偶数的点的个数分别应为 \(m\) 和 \(32-m\). 另一方面，在 \(T_1\) 中，横坐标可取 \(1\)，\(2\)，\(7\)，\(8\) 四个数，纵坐标可取 \(1\) 到 \(8\) 八个数. 对每个固定的横坐标，八个纵坐标中恰有四个与它同奇偶，四个与它异奇偶，因此横纵坐标之和为奇数的点共有 \(4\times 4=16\) 个，横纵坐标之和为偶数的点也共有 \(16\) 个，所以 \(m=16\). 再设
\(T_3=\{(x,y)\mid x\in A,\ y\in\{1,2,7,8\}\}\)，\(T_4=\{(x,y)\mid x\in A,\ y\in\{3,4,5,6\}\}\). 交换 \(x\) 与 \(y\) 的角色，按上面的论证可得 \(T_3=\{(x_k,y_k)\mid k=1,3,5,\dots,31,34,36,\dots,64\}\). 而前面对 \(T_1\) 也得到同样的指标集合，因此必有 \(T_1=T_3\). 但 \(T_1\) 与 \(T_3\) 的定义不同，不可能相等，矛盾. 所以 \(M\) 中所有元素都不能构成 \(k\) 列.
\end{enumerate}
\end{solution}
